%% ch02_mathematical_foundations.tex
%% Chapter 2: Mathematical Foundations of Model-Based Adaptive Protection
%% ============================================================
\chapter{Mathematical Foundations of Model-Based Adaptive Protection}
\label{ch:foundations}

\section{Introduction}
\label{sec:ch2_intro}

Every adaptive protection algorithm in this thesis rests on a common
mathematical structure: an \emph{inverse estimation problem}, in which
a physical model of the protected device is fitted to real-time measurements,
and the identified parameters are used to update relay thresholds. This
chapter establishes the theoretical foundations of that structure:
the general nonlinear least-squares formulation (\S\ref{sec:ch2_lsq}),
the Levenberg--Marquardt algorithm used as the solver
(\S\ref{sec:ch2_lm}), the composite confidence score that gates adaptation
decisions (\S\ref{sec:ch2_confidence}), and the bounded update law that
guarantees protection security (\S\ref{sec:ch2_update}).

These foundations are drawn upon directly by the generator OC
(Chapter~\ref{ch:generator_oc}) and transformer differential
(Chapter~\ref{ch:transformer_diff}) algorithms, and are reflected in
the EKF formulation of the digital twin (Chapter~\ref{ch:digital_twin}).

\section{Inverse Estimation Formulation}
\label{sec:ch2_lsq}

\subsection{General Structure}

Let $\mathbf{x}(t; \bm{\theta}) \in \mathbb{R}^n$ be the state vector
of a protected device (fault current, voltage, or impedance trajectory)
parameterised by an unknown vector $\bm{\theta} \in \mathbb{R}^p$.
Let $\mathbf{y}_k \in \mathbb{R}^n$ be a sequence of $K$ noisy measurements
at times $t_k$. The \emph{inverse estimation problem} is:

\begin{equation}
  \hat{\bm{\theta}} = \argmin_{\bm{\theta} \in \Theta}
    \sum_{k=1}^{K} w_k \left\|\mathbf{y}_k - \mathbf{x}(t_k;\bm{\theta})\right\|^2
  \label{eq:ch2_lsq}
\end{equation}

where $\Theta \subset \mathbb{R}^p$ is the feasible set defined by physical
bounds, $w_k > 0$ are sample weights (typically $w_k = 1$ for uniform
confidence within the measurement window), and $\|\cdot\|$ is the Euclidean
norm.

Writing $\mathbf{r}(\bm{\theta}) = \mathbf{y} - \mathbf{x}(\bm{\theta})$
as the residual vector (stacking all $K$ time steps), the problem becomes:
\begin{equation}
  \hat{\bm{\theta}} = \argmin_{\bm{\theta} \in \Theta}
    \frac{1}{2}\|\mathbf{r}(\bm{\theta})\|^2
  \label{eq:ch2_rss}
\end{equation}

When $\mathbf{x}(t;\bm{\theta})$ is nonlinear in $\bm{\theta}$ (as is the
case for all physical models in this thesis), Eq.~\eqref{eq:ch2_rss} is a
\emph{bounded nonlinear least-squares} problem.

\subsection{Jacobian and Sensitivity Analysis}

The first-order necessary condition $\nabla_{\bm{\theta}}(\frac{1}{2}\|\mathbf{r}\|^2) = 0$
gives the \emph{normal equations}:
\begin{equation}
  J(\bm{\theta})^\top \mathbf{r}(\bm{\theta}) = \mathbf{0}
  \label{eq:ch2_normal}
\end{equation}
where $J \in \mathbb{R}^{(nK) \times p}$ is the Jacobian:
\begin{equation}
  J_{ij} = \frac{\partial x_i}{\partial \theta_j}
  \label{eq:ch2_jacobian}
\end{equation}

The condition number $\kappa(J)$ measures the sensitivity of the solution
to measurement noise. When columns of $J$ are nearly linearly dependent
(ill-conditioning), small measurement perturbations cause large parameter
errors. Column-normalising $J$ before evaluating $\kappa$ removes
dimension effects:
\begin{equation}
  \kappa_n(J) = \kappa\left(\mathrm{diag}(\|J_{\cdot,1}\|,\ldots,\|J_{\cdot,p}\|)^{-1} J\right)
  \label{eq:ch2_kappan}
\end{equation}

\section{The Levenberg--Marquardt Algorithm}
\label{sec:ch2_lm}

\subsection{Derivation}

The Levenberg--Marquardt (LM) algorithm~\citep{Levenberg1944,Marquardt1963}
interpolates between steepest-descent (far from minimum) and
Gauss--Newton (near minimum) by adding a trust-region regulariser:

\begin{equation}
  \bm{\delta}^{(i)} = -(J^\top J + \mu^{(i)} I)^{-1} J^\top \mathbf{r}
  \label{eq:ch2_lm_step}
\end{equation}

\begin{equation}
  \bm{\theta}^{(i+1)} = \Pi_\Theta\left(\bm{\theta}^{(i)} + \bm{\delta}^{(i)}\right)
  \label{eq:ch2_lm_update}
\end{equation}

where $\Pi_\Theta(\cdot)$ projects onto the feasible set $\Theta$ (element-wise
clipping to physical bounds), and $\mu^{(i)} > 0$ is the damping parameter.

\subsection{Damping Parameter Adaptation}

The damping parameter is adapted by the gain ratio:
\begin{equation}
  \rho^{(i)} = \frac{\|\mathbf{r}(\bm{\theta}^{(i)})\|^2
                   - \|\mathbf{r}(\bm{\theta}^{(i+1)})\|^2}
                    {L(\bm{\theta}^{(i)}) - L(\bm{\theta}^{(i+1)})}
  \label{eq:ch2_gain_ratio}
\end{equation}
where $L(\bm{\theta})$ is the linearised model prediction:
$L(\bm{\theta}^{(i)}) - L(\bm{\theta}^{(i+1)}) = \bm{\delta}^\top(2J^\top \mathbf{r} - J^\top J\bm{\delta})$.

The update rule is:
\begin{equation}
  \mu^{(i+1)} = \begin{cases}
    \mu^{(i)} / 3  & \text{if } \rho^{(i)} > 0.75 \text{ (good step: reduce damping)}\\
    \mu^{(i)}      & \text{if } 0.25 \leq \rho^{(i)} \leq 0.75 \\
    \mu^{(i)} \times 10 & \text{if } \rho^{(i)} < 0.25 \text{ (poor step: increase damping)}
  \end{cases}
  \label{eq:ch2_mu_adapt}
\end{equation}

\subsection{Two-Pass Strategy}

For physical models where parameters separate into ``fast'' (identifiable
early in the measurement window) and ``slow'' (identifiable only at the tail
of the window) components, a two-pass approach improves convergence:

\begin{enumerate}
\item \textbf{Pass 1 (transient identification):} Fit the model over the
      early window $[t_0, t_0 + T_1]$ with the slow parameters held fixed at
      nominal values. Estimate the transient parameters $\bm{\theta}_\mathrm{fast}$.
\item \textbf{Pass 2 (steady-state identification):} Fit over the full
      window $[t_0, t_0 + T_2]$ with $\bm{\theta}_\mathrm{fast}$ fixed at
      Pass~1 values. Estimate the slow parameters $\bm{\theta}_\mathrm{slow}$.
\end{enumerate}

This strategy reduces the effective Jacobian condition number by decoupling
the transient and steady-state parameter subspaces, and is used in
the SAMBP-SyncOC algorithm (Chapter~\ref{ch:generator_oc}).

\begin{algorithm}[htbp]
\caption{Levenberg--Marquardt Inverse Estimator}
\label{alg:ch2_lm}
\begin{algorithmic}[1]
\Require Measurement window $\mathbf{y}$, model $\mathbf{x}(\cdot;\bm{\theta})$,
         initial guess $\bm{\theta}^{(0)}$, bounds $\Theta$,
         $\mu_0$, $\mu_\mathrm{max}$, $N_\mathrm{iter}$, $\varepsilon_r$
\Ensure  $\hat{\bm{\theta}}$, residual $R^2$, Jacobian $J$
\State $i \gets 0$,\; $\mu \gets \mu_0$
\While{$i < N_\mathrm{iter}$ and $\|\mathbf{r}^{(i)}\| > \varepsilon_r$}
  \State Compute $\mathbf{r}^{(i)} = \mathbf{y} - \mathbf{x}(\bm{\theta}^{(i)})$
  \State Compute Jacobian $J^{(i)}$ via finite differences
  \State Solve $(J^\top J + \mu I)\bm{\delta} = -J^\top \mathbf{r}^{(i)}$
  \State $\bm{\theta}^{\mathrm{cand}} \gets \Pi_\Theta(\bm{\theta}^{(i)} + \bm{\delta})$
  \State Compute gain ratio $\rho$ \Comment{Eq.~\eqref{eq:ch2_gain_ratio}}
  \If{$\rho > 0$}
    \State $\bm{\theta}^{(i+1)} \gets \bm{\theta}^{\mathrm{cand}}$
  \Else
    \State $\bm{\theta}^{(i+1)} \gets \bm{\theta}^{(i)}$ \Comment{reject step}
  \EndIf
  \State Update $\mu$ per Eq.~\eqref{eq:ch2_mu_adapt};\; $i \gets i+1$
\EndWhile
\State $\hat{\bm{\theta}} \gets \bm{\theta}^{(i)}$
\State Compute $R^2 = 1 - \|\mathbf{r}^{(i)}\|^2 / \|\mathbf{y} - \bar{\mathbf{y}}\|^2$
\end{algorithmic}
\end{algorithm}

\subsection{Convergence Properties}

Under standard regularity conditions~\citep{Nocedal2006}, the LM algorithm
satisfies:
\begin{equation}
  \|\bm{\theta}^{(i+1)} - \hat{\bm{\theta}}\| \leq C \|\bm{\theta}^{(i)} - \hat{\bm{\theta}}\|^{1+q}
  \label{eq:ch2_lm_convergence}
\end{equation}
with $q \in [1, 2)$ (superlinear convergence near the minimum, when $\mu \to 0$).
For the protection application, convergence to $R^2 > 0.999$ within 15--20
iterations has been demonstrated across all study cases (Chapter~\ref{ch:generator_oc}).

\section{Composite Confidence Score}
\label{sec:ch2_confidence}

After solving the inverse estimation problem, it is essential to know
\emph{how much to trust the result} before allowing it to update relay
settings. A wrong estimate with high apparent $R^2$ can arise from
overfitting, poor identifiability, or measurement noise. A four-component
composite confidence score $\gamma \in [0,1]$ is defined as:

\begin{equation}
  \gamma = w_1 \gamma_R + w_2 \gamma_J + w_3 \gamma_B + w_4 \gamma_P
  \label{eq:ch2_gamma}
\end{equation}

with $w_1 + w_2 + w_3 + w_4 = 1$ (default: $w = [0.35, 0.30, 0.20, 0.15]$).

\paragraph{Component $\gamma_R$ --- Normalised Residual Quality.}
\begin{equation}
  \gamma_R = \frac{1}{1 + \alpha_R (1-R^2)}
  \label{eq:ch2_gamma_r}
\end{equation}
where $\alpha_R = 100$ sharpens the score at $R^2 \approx 1$.
$\gamma_R \to 1$ when $R^2 \to 1$ (near-perfect fit).

\paragraph{Component $\gamma_J$ --- Jacobian Condition Number.}
\begin{equation}
  \gamma_J = \frac{1}{1 + \alpha_J \log_{10}(\kappa_n(J))}
  \label{eq:ch2_gamma_j}
\end{equation}
with $\alpha_J = 0.1$. $\gamma_J \to 1$ when $\kappa_n(J) \to 1$
(well-conditioned, identifiable problem); $\gamma_J \to 0$ as
$\kappa_n(J) \to \infty$ (unidentifiable).

\paragraph{Component $\gamma_B$ --- Bounds-Hit Fraction.}
\begin{equation}
  \gamma_B = 1 - \frac{N_\mathrm{bounds}}{p}
  \label{eq:ch2_gamma_b}
\end{equation}
where $N_\mathrm{bounds}$ is the number of parameters that converged to
a boundary of $\Theta$. Parameters hitting bounds indicate that the
optimum is physically infeasible or the model is misidentified.

\paragraph{Component $\gamma_P$ --- Physics Plausibility.}
$\gamma_P \in \{0, 0.5, 1.0\}$ is a rule-based score on the
physical plausibility of $\hat{\bm{\theta}}$. Each estimated parameter
is checked against physics-based secondary bounds:
\begin{equation}
  \gamma_P = 1 - \frac{N_\mathrm{implausible}}{p}
  \label{eq:ch2_gamma_p}
\end{equation}
where $N_\mathrm{implausible}$ is the number of parameters outside
physically expected ranges (e.g., $\hat{X}_d'' > \hat{X}_d'$,
$\hat{\tau}_{ac} < 0$).

\paragraph{Adaptation Gate.}
Adaptation of relay settings is authorised only when:
\begin{equation}
  \gamma \geq \gamma_\mathrm{th}
  \label{eq:ch2_gate}
\end{equation}
with default $\gamma_\mathrm{th} = 0.70$.
This threshold was chosen to balance responsiveness (adapt quickly to
real changes) against security (reject spurious estimates).

\section{Bounded Adaptive Update Law}
\label{sec:ch2_update}

\subsection{Motivation}

If the inverse estimator produces an erroneous $\hat{\bm{\theta}}$ that
passes the confidence gate (false positive), the relay settings must not
move to an unsafe operating point. The bounded update law imposes three
levels of safety:

\begin{enumerate}
\item \textbf{Rate limiting:} Settings change at most $\delta_\mathrm{max}$ per
      update cycle, preventing step changes.
\item \textbf{Hard bounds:} Settings are clamped to $[\theta_\mathrm{min},
      \theta_\mathrm{max}]$ from commissioning specifications.
\item \textbf{Security floor:} Settings may never drop below the minimum
      value required to detect the worst-case bolted fault with a 20\%
      margin.
\end{enumerate}

\subsection{Update Equation}

Let $\xi_k$ be the $k$-th relay setting (e.g., pickup $I_p$ or
time-multiplier $\TMS$), and $\hat{\xi}_k(\hat{\bm{\theta}})$ the
value recommended by the identified model. The bounded update is:

\begin{equation}
  \xi_k^{(n+1)} = \mathrm{clip}\!\left(
    \xi_k^{(n)} + \delta_\mathrm{max}\,
    \mathrm{sgn}(\hat{\xi}_k - \xi_k^{(n)})
    \cdot \min\left(1,\,
      \frac{|\hat{\xi}_k - \xi_k^{(n)}|}{\delta_\mathrm{max}}
    \right),\;
    \xi_\mathrm{min},\; \xi_\mathrm{max}
  \right)
  \label{eq:ch2_update}
\end{equation}

For the IEC Very-Inverse relay:
\begin{align}
  I_p^{(n+1)} &= \mathrm{clip}\!\left(I_p^{(n)} + \Delta I_p,\,
    1.05 I_\mathrm{load}^{\max},\, 0.80 I_\mathrm{fault}^{\min}\right)
  \label{eq:ch2_ip_update} \\
  \TMS^{(n+1)} &= \mathrm{clip}\!\left(\TMS^{(n)} + \Delta \TMS,\,
    \TMS_\mathrm{min},\, \TMS_\mathrm{max}\right)
  \label{eq:ch2_tms_update}
\end{align}

\subsection{Security Proof}

\begin{proposition}[Bounded update preserves trip security]
\label{prop:ch2_security}
If $\xi_k^{(0)}$ is in $[\xi_\mathrm{min}, \xi_\mathrm{max}]$ and all
$\hat{\xi}_k$ are in $\mathbb{R}$, then for all $n \geq 0$:
$\xi_k^{(n)} \in [\xi_\mathrm{min}, \xi_\mathrm{max}]$.
\end{proposition}

\begin{proof}
The update Eq.~\eqref{eq:ch2_update} applies $\mathrm{clip}(\cdot,
\xi_\mathrm{min}, \xi_\mathrm{max})$ at every step. Since $\mathrm{clip}$
is idempotent and bounds-preserving, the invariant holds by induction.
\end{proof}

\section{Applicability Map}
\label{sec:ch2_map}

Table~\ref{tab:ch2_applicability} summarises how the inverse estimation
framework specialises to each protection element in this thesis.
The mathematical structure (Eq.~\eqref{eq:ch2_lsq}--\eqref{eq:ch2_update})
is identical; only the physical model $\mathbf{x}(t;\bm{\theta})$ and
the parameter vector $\bm{\theta}$ change.

\begin{table}[htbp]
\caption{Inverse estimation framework specialisation by protection element}
\label{tab:ch2_applicability}
\small
\begin{tabularx}{\textwidth}{@{}lp{2.5cm}p{3.5cm}p{2.5cm}p{2cm}@{}}
\toprule
Element & Model & Parameters $\bm{\theta}$ & Adapted setting & Status \\
\midrule
87L line diff & IBR differential current & $k_\mathrm{ibr}$, $\phi_B$ & Bias slope $k_m^*$ & [P] \\
Gen. OC (SG) & 3-component fault current & $\{I_{ss},I_{sub},\tau_{ac},\tau_{dc},\phi_a\}$ & $I_p$, $\TMS$ & [P] \\
Gen. OC (DFIG) & Piecewise crowbar model & $\{I_{0+},\tau_r,I_{crowd}\}$ & $I_p$ & [F] \\
87T transformer & Inrush vs.\ fault model & $\{i_2/i_1,\phi_\mathrm{hm}\}$ & SPRT threshold & [F] \\
EKF digital twin & IBR Thevenin model & $\kibr$, $\Zline$, $\alpha$ & 87L bias $k_m^*$ & [P] \\
\bottomrule
\end{tabularx}
\end{table}

\section{Relationship to Kalman Filtering}
\label{sec:ch2_ekf_relation}

The EKF used in the digital twin (Chapter~\ref{ch:digital_twin}) is
formally a \emph{recursive} solution to the inverse estimation problem
with Gaussian process and measurement noise. Writing the state model as:
\begin{align}
  \mathbf{x}_k &= f(\mathbf{x}_{k-1}) + \mathbf{w}_k, \quad
    \mathbf{w}_k \sim \mathcal{N}(\mathbf{0}, Q) \\
  \mathbf{y}_k &= h(\mathbf{x}_k) + \mathbf{v}_k, \quad
    \mathbf{v}_k \sim \mathcal{N}(\mathbf{0}, R)
\end{align}
the EKF update step minimises the same objective as
Eq.~\eqref{eq:ch2_lsq} over a sliding window, with the prediction
step providing the Tikhonov regularisation that plays the role of the
LM damping term $\mu$. This formal connection means the confidence
score framework of \S\ref{sec:ch2_confidence} applies directly to
evaluate the quality of EKF state estimates.
